\problemname{Canteen Rush}
At 12:00 each day, the lecture ends for a lot of students at the ÆT-University. The ravenous students rush towards the canteens of the university to fulfill their hunger. 
However, there is a problem. The doors of the university are incredibly shabby, definitely not up to code and they are one way \(!?\). They can only handle a certain number of people 
rushing through during lunch break before they break down for the day, trapping the poor students who didn’t make it through.
\\\\
Students start at different auditoriums around the campus, and must each make their way to one of the canteens. The university is a big building with numerous connecting rooms, 
each named uniquely . Luckily, the university wasn't completely heartless when designing the campus -- each auditorium has at least one path to a canteen. 
Your job is to find the maximum number of students who can theoretically make it to lunch during these lunch rush hours.

\section*{Input}
\begin{itemize}
    \item One line with an integer $A$ ($1 \leq A \le wip$), the number of auditoriums.
    \item $A$ lines, each containing the name of the auditorium followed by an integer $p$, the number of people in the auditorium. The name consists of at most 20 lowercase characters, including digits and underscores.
    \item One line with an integer $C$ ($1 <\leq C \le wip$), the number of canteens
    \item $C$ lines, each line containing the name of a canteen. The name consists of at most 20 lowercase characters, including digits and underscores.
    \item One line with an integer $D$ ($1 \leq D \le wip$), the number of doors connecting the rooms.
    \item $D$ lines containing room names $u$ and $v$, denoting that room u and v are connected by a door, and an integer $c$, denoting the capacity of the door. The names consists of at most 20 lowercase characters, including digits and underscores.


\end{itemize}

\section*{Output}
Output one integer, the maximum number of people who make it to a canteen.
